\documentclass[twoside]{article}
\usepackage{aistats2026}
\usepackage[round]{natbib}
\renewcommand{\bibname}{References}
\renewcommand{\bibsection}{\subsubsection*{\bibname}}
\bibliographystyle{apalike}
\usepackage[hyphens]{url}
\usepackage{graphicx}
\urlstyle{rm}
\def\UrlFont{\rm}
\usepackage{amsmath,amssymb,amsthm,mathtools,bm}
\usepackage{booktabs}
\usepackage{multirow}
\usepackage{tikz}
\usetikzlibrary{arrows.meta,positioning,decorations.pathreplacing}
\special{papersize = 8.5in, 11in}
\setlength{\pdfpageheight}{11in}
\setlength{\pdfpagewidth}{8.5in}
\newtheorem{proposition}{Proposition}
\newtheorem{lemma}{Lemma}
\newtheorem{remark}{Remark}
\newcommand{\E}{\mathbb E}
\newcommand{\R}{\mathbb R}
\newcommand{\M}{\mathcal M}
\newcommand{\D}{\mathcal D}
\newcommand{\A}{\mathcal A}
\newcommand{\sg}{\operatorname{sg}}
\DeclareMathOperator*{\argmin}{arg\,min}

\begin{document}
\twocolumn[
\aistatstitle{PART: Proximal Actor Refinement through Target Routing in Offline Reinforcement Learning}
\aistatsauthor{Anonymous Authors}
\aistatsaddress{Anonymous Institution}
]

\begin{abstract}
The squared TD3+BC anchor admits a JKO-style Wasserstein--2 interpretation:
under a fixed critic and exact near-identity updates, splitting nominal time
$T$ across $K$ hops discretizes one local gradient flow. We study the
finite-$T$ regime where this picture fails. Proximal actor refinement through
target routing (PART) routes only the first of $K$ persistent actors into TD
targets and deploys the last. Across nine D4RL tasks, fourteen budgets, and
four seeds, the procedures are close at small $T$, but at $T=4,7$ TD3+BC falls
to $45.93,30.36$ while PART-P4 remains at $82.06,81.79$. Over $T\ge4$,
increasing $K$ from one to four raises the task-equal mean from $25.38$ to
$63.97$ and cuts score-$<20$ cells from $34/63$ to $8/63$; both disjoint seed
pairs reproduce the ordering. Matched audits show that P3/P4 reduce
first-branch action and target-value perturbation without typical final-policy
contraction. A two-actor P3 control reaches the same seed-mean collapse count
without sequential re-centering. In a completed 90-pair P4 comparison, PART
minus control has mean $-2.31$ and interval $[-7.25,2.41]$; the predeclared
outcome is unresolved. Thus PART shifts the finite-step stability boundary;
matched evidence is consistent with exposure--reach separation, while the
chain's incremental return benefit is not established.
\end{abstract}

\section{INTRODUCTION}
\label{sec:introduction}

Behavior-regularized offline RL admits a Jordan--Kinderlehrer--Otto (JKO) minimizing-movement view \citep{jordan1998variational}. For deterministic policies, the squared TD3+BC anchor is a state-averaged Wasserstein--2 movement cost and its normalized coefficient defines nominal time $T=\alpha/2$. With a fixed critic, exact optimization, inactive projection, and $T\to0$, explicit and implicit updates follow the same local action-space flow to first order.

Our question begins where this local picture breaks. Offline RL fits the critic to policy-generated target actions, so increasing $T$ changes the bootstrapping distribution as well as a numerical step. Projection activates, neural actors amortize rather than solve subproblems, persistent actors retain distinct histories, and the critic moves. TD3+BC further uses a Polyak copy of the deployed actor in Bellman targets \citep{fujimoto2021minimalist}, coupling deployment reach to critic-facing movement and enabling critic error to feed back into policy improvement.

Proximal actor refinement through target routing (PART) allocates $T$ across $K$ sequential actors. The first is data anchored, later actors are anchored to stop-gradient predecessor outputs, a Polyak target copy of the first supplies TD-target actions, and the last is deployed. Let $D_{\mathrm{critic}}$ and $D_{\mathrm{final}}$ denote first- and final-branch displacement. A one-hop actor identifies them; PART separates them. PART is one structured realization of this target/deployment separation, not evidence by definition that sequential re-centering is necessary. We write PART-P$k$ (P$k$ for short) for the proximal $K=k$ condition and PART-L$k$ (L$k$) for its projected-linearized analogue. Released artifact paths retain the historical Budgeted Actor Refinement (\texttt{BAR}) identifier; the paper uses PART throughout.

This yields three interpretive regimes. The \emph{local regime} is the asymptotic fixed-critic limit. In the finite-step \emph{transition regime}, a single large actor step can move the Bellman target beyond a critic-valid neighborhood, whereas routing a smaller first hop into backups leaves later branches for deployment refinement. In the \emph{large-budget stress regime}, projection, amortized fitting, critic drift, and optimizer history dominate; depth can shift but not remove the failure boundary. These are descriptive empirical regimes, not universal theorem boundaries.

The phase diagram follows this pattern. At $T=.05,.1,.2$, depth has no consistent return ordering. TD3+BC then falls to $67.41,45.93,30.36$ at $T=2.5,4,7$, while P4 gives $81.25,82.06,81.79$. Across $T\ge4$, means rise $25.38\to41.20\to51.92\to63.97$ for $K=1,2,3,4$ and low-return cells fall $34\to22\to16\to8$. Matched P3/P4 audits find lower first-branch displacement in 88/90 and 89/90 cells and lower common-critic target-value perturbation in 86/90 and 88/90 cells, without typical final-policy contraction. A P3 two-actor control reproduces the aggregate regime without sequential re-centering; the completed P4 comparison is unresolved. Thus policy separation is the simplest compatible mechanism, while the sequential chain's incremental return benefit is not established.

Our contributions are as follows.
\begin{itemize}
    \item We formulate PART for TD3+BC as a re-centered actor composition that assigns target-policy displacement and final-policy reach to distinct branches.
    \item We connect the anchored update to an ideal Wasserstein/JKO flow,
    derive finite-step critic-error and target-perturbation statements, and
    identify where projection, ReLU boundaries, and live optimization end the
    small-step interpretation.
    \item Four-seed phase diagrams and matched audits show that multi-hop PART
    becomes useful as TD3+BC exits its small-budget plateau: it reduces the
    target-coupled first movement, preserves deployment reach, and shifts the
    large-budget collapse boundary, chiefly through lower-tail rescue.
\end{itemize}

\section{RELATED WORK}
\label{sec:related}

\paragraph{Behavior-regularized offline actor--critic learning.}
Offline actor--critic methods must control value extrapolation when a learned
policy selects actions outside the fixed dataset. Prior work studies this
failure through constrained policy improvement, conservative value learning,
and restrictions on actions used in Bellman backups
\citep{fujimoto2019offpolicy,kumar2019bear}. TD3+BC provides a particularly
minimal control: a deterministic actor maximizes a normalized critic while a
squared cloning term anchors its actions to the dataset
\citep{fujimoto2021minimalist}. We keep this critic update and actor family
fixed and study which refined actor branch supplies the TD target.
One-step offline RL instead limits repeated policy optimization against an
off-policy value estimate \citep{brandfonbrener2021offline}. PART retains
iterative actor--critic training but distributes a nominal improvement
coefficient across a routed actor chain.

\paragraph{Policy geometry and proximal refinement.}
Trust-region policy optimization is commonly derived from information geometry
\citep{amari2000methods,schulman2015trust}. Optimal-transport geometry gives a
different notion of movement: policy optimization can be expressed as a
Wasserstein gradient flow, with related natural-gradient and trust-region
constructions for continuous control
\citep{zhang2018policy,moskovitz2020efficient,terpin2022trust,song2023metric}.
Wasserstein Policy Optimization projects Wasserstein policy flow into a
parametric actor class using value-function action gradients
\citep{pfau2025wpo}.
For deterministic policies, the statewise Wasserstein--2 cost reduces to
squared action displacement. We use that identity to interpret the TD3+BC
anchor as a proximal movement cost, while separating the finite-step
variational statement from its small-step flow limit.

\paragraph{Target networks and actor--critic coupling.}
Bootstrapping error can feed critic extrapolation back into actor updates
\citep{fujimoto2019offpolicy,kumar2019bear}. Polyak averaging delays this
feedback but ordinarily copies the same actor branch that will be deployed.
MCEP instead pairs a conservative target policy for value learning with a
mildly constrained evaluation policy excluded from Bellman backups
\citep{xu2024mcep}. PART shares this routing principle but uses a persistent,
predecessor-recentered chain and varies depth under a fixed nominal budget;
MCEP keeps both actors directly data-anchored with separately chosen strengths
and is not a $K=2$ instance of PART. Our contribution concerns this depth-indexed
routing operator rather than a new critic objective.

\section{JKO VIEW OF TD3+BC}
\label{sec:background}

We consider a discounted Markov decision process
$\langle\mathcal S,\mathcal A,P,r,\gamma\rangle$ with continuous actions and
a fixed offline dataset $\mathcal D=\{(s_i,a_i,r_i,s_i')\}_{i=1}^N$. A clipped
double critic is trained from targets
\begin{equation}
y_{\rho}=r+\gamma\min_j Q_j^-(s',\rho(s')),
\label{eq:td_target_background}
\end{equation}
where $\rho$ is a target policy. When $\rho(s')$ leaves dataset support,
critic error enters the regression target and can be amplified by later actor
updates \citep{fujimoto2019offpolicy,kumar2019bear}.

For a deterministic actor $\mu_\theta$, normalized TD3+BC minimizes
\begin{equation}
\mathcal L_{\rm actor}(\theta)
=-\alpha\,\mathbb E_s\bar Q(s,\mu_\theta(s))
+\mathbb E_{(s,a)\sim\mathcal D}\frac{\|\mu_\theta(s)-a\|_2^2}{d},
\label{eq:td3bc_background}
\end{equation}
where $\bar Q=\widehat Q/C$ and $C>0$ is a stop-gradient critic scale
\citep{fujimoto2021minimalist}. Ordinarily, a Polyak copy of this actor enters
Eq.~\eqref{eq:td_target_background}, while the online actor is deployed.
Target-policy movement and deployment reach therefore follow the same branch.

A metric proximal update with step $h>0$ takes the form
\begin{equation}
\pi^+\in\argmin_\pi\left\{\widehat{\mathcal E}(\pi)
+\frac{1}{2h}d_{\mathcal M}^2(\pi,\pi^{\rm ref})\right\}.
\label{eq:generic_prox}
\end{equation}
This minimizing-movement construction is associated with metric gradient
flows \citep{ambrosio2005gradient,otto2001geometry}. For deterministic
state-conditioned policies,
\begin{equation}
W_2^2(\delta_{\mu(s)},\delta_{\nu(s)})=\|\mu(s)-\nu(s)\|_2^2.
\label{eq:dirac_w2_background}
\end{equation}
Thus the TD3+BC anchor is a Wasserstein--2 movement cost up to fixed scaling.
We distinguish its finite-step variational meaning from the small-step,
fixed-critic gradient-flow limit and from large-budget coupled training.

\section{FROM LOCAL FLOW TO PART}
\label{sec:method}

PART maintains $K$ independently initialized, persistent deterministic
actors with separate optimizer states. At each delayed update they take one
Adam step in chain order on the same minibatch; Figure~\ref{fig:algorithm} gives
the exact live ordering. The online critic is fixed within the actor chain; we
denote its first Q-function by $\widehat Q$. We first derive the proximal
operator motivating these updates and then state its live-algorithm boundary.
\subsection{Offline actor updates as proximal improvement}

Let $\rho_{\mathcal D}$ be the empirical state distribution and consider a
deterministic actor $\mu:\mathcal S\rightarrow\mathcal A\subset\mathbb R^d$.
At one training iteration, the critic parameters are held fixed while the actor
chain is updated. We write the current critic as $\widehat Q_t$, use
$\epsilon=10^{-6}$, and define the implementation-level proximal scales by
\begin{equation}
\begin{aligned}
 C_{1,t}^{\rm P}
 &=\sg\!\left(\mathbb E_s
 |\widehat Q_t(s,\mu_{1,t}^{\rm pre}(s))|+\epsilon\right),\\
 C_{k,t}^{\rm P}
 &=\sg\!\left(\mathbb E_s
 |\widehat Q_t(s,\sg\mu_{k-1,t}^{+}(s))|+\epsilon\right),\quad k\ge2.
\end{aligned}
\label{eq:qnorm}
\end{equation}
The normalized critic is
$\bar Q_{k,t}(s,a)=\widehat Q_t(s,a)/C_{k,t}^{\rm P}$.
Here \emph{pre} and the superscript $+$ denote actions before and after the
current actor update, and $\sg$ denotes stop-gradient. We suppress $t$ and the
superscript when discussing the ideal operator. A fixed positive $C$ only
rescales gradient-flow time; recomputing $C_k$ changes the finite-step schedule.
Our finite-step statements condition on the realized scales.

For deterministic policies and a state distribution $\rho$, define the
action-averaged Wasserstein metric
\begin{equation}
 \mathsf d_{d,\rho}^2(\mu,\nu)
 :=\mathbb E_{s\sim\rho}
 \left[\frac{1}{d}\|\mu(s)-\nu(s)\|_2^2\right].
 \label{eq:mean_metric}
\end{equation}
The squared quantity is $d^{-1}$ times the state-averaged squared $W_2$
distance between state-conditioned Dirac policies. Its square root
$\mathsf d_{d,\rho}$ is a pseudometric on pointwise policy maps and a metric
after identifying maps equal $\rho$-almost everywhere. We abbreviate
$\mathsf d_d=\mathsf d_{d,\rho_{\mathcal D}}$ on the actor-training state
distribution.
This is an $L_2(\rho)$ geometry on conditional Dirac policy maps with the
state marginal held fixed, not Wasserstein transport between policy-induced
state--action occupancy measures.
The population counterpart of squared behavior cloning also admits a
deterministic reference. Let $P_{\!B}$ be the data-generating behavior
distribution, let $\rho_B$ be its state marginal, and define
$b(s)=\mathbb E_{P_{\!B}}[a\mid s]$. Then
\begin{equation}
 \mathbb E_{(s,a)\sim P_{\!B}}\frac{1}{d}\|\mu(s)-a\|_2^2
 =\mathsf d_{d,\rho_B}^2(\mu,b)+\mathrm{const},
 \label{eq:bc_decomposition}
\end{equation}
where the constant is the conditional action variance and does not depend on
$\mu$. The empirical actor loss is a Monte Carlo estimate of the left-hand
side, so samplewise squared anchoring targets the conditional-mean behavior map
at the population level.

Given a reference policy $\nu$, step $h>0$, and fixed normalized critic
$\bar Q$, define the proximal refinement operator by
\begin{equation}
\mathcal P_{h}^{\bar Q}(\nu)
\in\argmin_{\mu\in\mathcal M}
\left\{-\mathbb E_s\bar Q(s,\mu(s))
+\frac{1}{2h}\mathsf d_d^2(\mu,\nu)\right\}.
\label{eq:prox_operator}
\end{equation}
At hop $k$, setting $\bar Q=\bar Q_k$ and multiplying by $2h$ gives the
practical actor objective
\begin{equation}
\mathcal L_{h,k}(\mu\mid\nu)
=-2h\,\mathbb E_s\bar Q_k(s,\mu(s))
+\mathsf d_d^2(\mu,\nu).
\label{eq:actorloss}
\end{equation}

\begin{proposition}[Behavior anchoring is Wasserstein proximal improvement]
\label{prop:prox_equiv}
For deterministic policies, the normalized anchored objective
\begin{equation*}
 -\alpha\,\mathbb E_s\bar Q(s,\mu(s))
 +\mathsf d_d^2(\mu,\nu)
\end{equation*}
is Eq.~\eqref{eq:prox_operator}, up to multiplication by a positive constant, with
$h=\alpha/2$. Equivalently, the raw-critic objective
$-\lambda\mathbb E_s\widehat Q(s,\mu(s))+\mathsf d_d^2(\mu,\nu)$ has
$h=\lambda C/2$ when $C$ is fixed. The population counterpart of squared
dataset anchoring uses $\nu=b$ by Eq.~\eqref{eq:bc_decomposition}.
\end{proposition}

The proposition identifies an update operator relative to the \emph{estimated}
critic; true-return improvement lies outside the scope of this certificate.

\subsection{Re-centered composition and actor--critic dataflow}

For population analysis, let $\mu_0:=b$ denote the conditional-mean behavior
map and let $h_{1:K}$ be positive step sizes. The idealized PART chain is
\begin{equation}
 \mu_k\in\mathcal P_{h_k}^{\bar Q_k}
 (\operatorname{sg}\mu_{k-1}),
 \qquad k=1,\ldots,K.
 \label{eq:bar}
\end{equation}
The implementation approximates each subproblem with one Adam update.
Its first actor $\mu_1$ uses sampled dataset actions as its reference; later
actors $\mu_2,\ldots,\mu_K$ use the preceding online actor under the same
current critic. Only the Polyak copy of $\mu_1$ enters the TD target, while
$\mu_K$ supplies the evaluated action. Because the later persistent actors are
independently initialized, their anchoring gradient need not vanish even at
$h=0$ when predecessor and current actor differ. The ideal proximal map
therefore motivates, but does not certify, the live one-Adam-step chain.

\begin{figure}[t]
\centering
\small
\fbox{\begin{minipage}{.94\columnwidth}
\textbf{One PART training iteration.} The $K$ actors are persistent,
independently initialized, and have separate Adam states.
\vspace{2pt}
\begin{tabular}{@{}rp{.82\linewidth}@{}}
1. & Form the Bellman target with $\mu_1^-$ and update the online critic.\\
2. & On a delayed step, update $\mu_1$ once toward sampled dataset actions.\\
3. & Polyak-update $\mu_1^-$ and the target critics with rate $.005$.\\
4. & For $k=2{:}K$, use the updated $\sg\mu_{k-1}^{+}$ as reference and update $\mu_k$ once on the same minibatch, holding the online critic fixed.\\
5. & Use $\mu_1^-$ in future Bellman targets and deploy $\mu_K$.\\
\end{tabular}
\end{minipage}}
\caption{Actual update and routing order. Each hop is one optimizer step, not
an exact proximal solve.}
\label{fig:algorithm}
\end{figure}

Our fixed-budget schedule uses the implementation parameter
$T=\tau=\alpha/2$, sets $h_k=T/K$, and holds $T$ fixed. Thus $K=1$ recovers
the TD3+BC coefficient $\alpha=2T$; $T$ is neither a compute budget nor a
realized distance constraint. Increasing $K$ jointly changes local coefficient,
re-centering frequency, actor count and compute, normalization trajectory, and
target/deployment dataflow. Each $K$ is therefore a complete randomized
procedure, and its seed average estimates marginal end-to-end performance.
Identical critic initialization across depths is a common-random-number
variance-reduction design, not a validity condition.

On the next-state distribution $\rho_{\mathcal D}'$ used by the TD target,
define the action-averaged population displacements
\begin{equation}
\begin{aligned}
D_{\rm critic}&:=\mathsf d_{d,\rho_{\mathcal D}'}^2(\mu_1,\mu_0),\\
D_{\rm final}&:=\mathsf d_{d,\rho_{\mathcal D}'}^2(\mu_K,\mu_0).
\end{aligned}
\label{eq:displacements}
\end{equation}
A one-step actor has $D_{\rm critic}=D_{\rm final}$; PART assigns the two
population quantities to distinct policy branches. On fixed validation batches, we measure the sample-anchored online branches
and the deterministic next-state Polyak target branch:
\begin{equation}
\begin{aligned}
\widehat D_{\rm critic}&=\mathbb E_i\|\mu_1(s_i)-a_i\|_2^2/d,\\
\widehat D_{\rm final}&=\mathbb E_i\|\mu_K(s_i)-a_i\|_2^2/d,\\
\widehat D_{\rm target}&=\mathbb E_i\|\mu_1^-(s_i')-a_{i+1}\|_2^2/d.
\end{aligned}
\label{eq:direct_target_displacement}
\end{equation}
The direct audit retains only bootstrapped transition pairs; a smoothed version
uses common TD3 target-noise draws. Sample anchoring adds a common
conditional-action-variance floor in population differences. These are
movement measurements, not critic-reliability estimates.

\subsection{Finite-step statements}

For a re-centered hop where $\mu_{k-1}$ is feasible, suppose the realized
actor satisfies
$\mathcal L_{h_k,k}(\mu_k\mid\mu_{k-1})\le
\mathcal L_{h_k,k}(\mu_{k-1}\mid\mu_{k-1})+\varepsilon_k$.
Rearranging the practical objective gives the conditional estimated-critic
margin
\begin{equation}
\begin{split}
 m_k^{\widehat Q}:={}&\frac{2h_k}{C_k}\mathbb E_s
 [\widehat Q(s,\mu_k(s))-\widehat Q(s,\mu_{k-1}(s))]\\
 &-\mathsf d_d^2(\mu_k,\mu_{k-1})\ge-\varepsilon_k.
\end{split}
\label{eq:margin}
\end{equation}
For $\widehat Q=Q^\dagger+e$, this transfers to a trusted objective as
\begin{equation}
 m_k^{Q^\dagger}=m_k^{\widehat Q}
 -\frac{2h_k}{C_k}\mathbb E_s[e(s,\mu_k(s))-e(s,\mu_{k-1}(s))].
\label{eq:critic_error_margin}
\end{equation}
The first sample-anchored hop lacks a feasible deterministic dataset-action
comparator, and one Adam step supplies no bound on the unmeasured
$\varepsilon_k$. These statements therefore diagnose a realized learned
objective; they guarantee neither exact proximal solution nor monotone return.
The Supplement gives the derivation and bounded-error corollary.

\begin{proposition}[First-hop TD-target perturbation bound]
\label{prop:exposure}
For clarity, consider a noiseless TD target using $\mu_1$,
$y_{\mu_1}=r+\gamma\min_jQ_j^-(s',\mu_1(s'))$, and suppose each target critic
$Q_j^-(s',\cdot)$ is $L(s')$-Lipschitz. For any reference $\mu_0$,
\begin{equation}
\begin{split}
\|y_{\mu_1}-y_{\mu_0}\|_{L_2(\mathcal D)}
\leq\gamma\Big(&\mathbb E_{s'}L(s')^2\\
&\cdot\|\mu_1(s')-\mu_0(s')\|_2^2\Big)^{1/2}.
\end{split}
\label{eq:exposure_bound}
\end{equation}
If $L(s')\leq\bar L$, the right-hand side is at most
$\gamma\bar L\sqrt{dD_{\rm critic}}$. Conditional on $\mu_1$, later actors do
not enter this one-step target perturbation.
\end{proposition}

Proposition~\ref{prop:exposure} bounds a one-step value-target difference, not
the stability of the learned critic. Its architectural content is that later
refinement can alter $D_{\rm final}$ without directly entering this bound.
With a Polyak target network, $\mu_1$ is replaced by its target copy. A
same-noise comparison under TD3 target smoothing obeys the same bound because
action clipping is nonexpansive. The current-state
$\widehat D_{\rm critic}$ omits Polyak lag and target noise, whereas
$\widehat D_{\rm target}$ measures the Polyak actor directly and has a
smoothed counterpart. Neither displacement measures the action Lipschitz
factor or the resulting change in critic value.

The separation introduces a complementary evaluation-policy mismatch: critic
targets follow the $\mu_1$ branch, whereas evaluation follows $\mu_K$.
Increasing $D_{\rm final}$ at fixed $D_{\rm critic}$ can therefore make the
critic less representative of the evaluated policy. We interpret PART as
exposing a measurable target/deployment trade-off; its empirical return
captures the net effect of target movement, evaluation mismatch, actor compute,
and the coupled optimization trajectory.

\subsection{Small-budget flow regime and its boundary}

With a fixed critic and common positive scale $C$, inactive projection, and an
exact statewise solve, the linearized step is $\beta g$ and the ideal proximal
step is $\beta g+\beta^2Hg+O(\beta^3)$; both are first order. Repeated
composition of one consistent near-identity map has fixed-$K$, $T\to0$ flow
bias of order $T^2/K$ (Supplement). Persistent independently initialized
actors need not satisfy that premise. This limit does not explain the coupled
large-$T$ chain. For a ReLU critic, the explicit update and the same-region
stationary proximal branch coincide before the first activation or box
boundary. In a prelocked
fresh-index audit of previously studied Hopper/Walker $T=1$ critics, task-equal
residence over 6144 eligible interior anchors is $.954,.911,.842,.748$ at probe
horizons $.025,.05,.1,.2$ (independent-implementation verifier pass); the
$.2$-horizon task range is $.310$--$.980$. This sharply delimits the local
identity
without asserting a return effect, global certificate, or learned $1/K$ slope.

\section{EXPERIMENTS}
\label{sec:experiments}

We evaluate PART on the nine D4RL locomotion tasks formed by HalfCheetah,
Hopper, and Walker2d with medium, medium-replay, and expert datasets
\citep{fu2020d4rl}. These tasks directly exercise deterministic continuous
control and provide distinct behavior-data regimes for a target-routing stress
test. Base implementations follow CORL configurations
\citep{tarasov2022corl}; scores are D4RL normalized returns. Additional
per-task results and diagnostic definitions are in
the Supplementary Material.

\paragraph{Fixed-budget protocol and estimand.}
We hold a total nominal budget $T$ fixed, split it across $K$ actor hops as
$h=T/K$, and sweep
\begin{equation*}
\begin{split}
T\in\{&0.05,0.1,0.2,0.4,0.7,1.5,2.5,\\
      &4,7,10,12,14,17,20\}.
\end{split}
\end{equation*}
over all nine tasks. The proximal and projected-linearized grids use seeds 0--3, for
$9\times14\times4=504$ runs per depth.
Proximal seeds 0--1 and 2--3 were executed on separate hosts and form disjoint
replication pairs. The grid is a high-resolution stability phase diagram rather
than a final benchmark alone. For $K=1$, Proposition~\ref{prop:prox_equiv}
maps $T$ to the normalized TD3+BC coefficient $\alpha=2T$.

We compare TD3+BC ($K=1$), proximal PART with $K\in\{2,3,4\}$ (PART-Prox), and an
action-metric-matched projected linearized target with $K\in\{2,3\}$ (PART-Lin):
\begin{equation}
\begin{aligned}
 a_k^{\rm tar}
 &=\Pi_{\mathcal A}\!\left[
 \nu_k+d\frac{T}{K}\frac{\nabla_a\widehat Q(s,\nu_k)}{C_{k,t}^{\rm L}}\right],\\
 \nu_1&=a,\qquad
 \nu_k=\sg\mu_{k-1,t}^{+}(s)\quad(k\ge2).
\end{aligned}
 \label{eq:linearized_target}
\end{equation}
At each hop, PART-Lin detaches the reference, recomputes the $\widehat Q$
gradient and scale, clips the target to the action box, and takes one
mean-squared-error optimizer step with
the corresponding persistent actor. Its first-hop scale is
\begin{equation*}
C_{1,t}^{\rm L}=\sg\!\left(
\mathbb E_{(s,a)\sim\mathcal D}|\widehat Q_t(s,a)|+\epsilon\right),
\end{equation*}
whereas $C_{k,t}^{\rm L}=C_{k,t}^{\rm P}$ for $k\ge2$. Thus the factor $d$
matches the action metric, but the first-hop normalization point differs from
PART-Prox; the comparison does not match the complete scale schedule. It uses the same
first-actor Bellman-target and final-actor evaluation routing. All conditions match
datasets, critic updates, and the $10^6$-update horizon; every headline score is
the ten-episode D4RL normalized mean at the final checkpoint. Intermediate
evaluation cadence varies by run family and does not enter this estimand. A
$K$-hop condition performs $K$ actor optimizer steps per delayed update, so
depth comparisons are between complete training procedures. The end-to-end
intervention changes local coefficient, re-centering frequency, actor compute,
normalization trajectory, routing, and optimization history.

We test whether this intervention reduces large-budget low-return outcomes,
persists under the linearized realization, separates first-hop exposure from
final-policy reach, and exposes critic-validity signatures of failure.

\subsection{The advantage emerges beyond the local regime}

Figure~\ref{fig:stability} and Table~\ref{tab:stability} reveal three empirical
regions. On $T=.05,.1,.2$, which overlaps the independently probed local
horizons, depth gives no consistent return ordering. TD3+BC peaks at $T=1.5$
and then leaves the high-return plateau: its mean is $67.41,45.93,30.36$ at
$T=2.5,4,7$, whereas P4 gives $81.25,82.06,81.79$. At $T\ge10$ every method
eventually degrades, but the proximal tail remains depth ordered. These are
descriptive regime labels, not predeclared bins. The residence audit and the
$T=20$ projection diagnostic bracket, but do not directly locate, where the
local assumptions fail.

The high-budget estimand averages seeds and the seven $T\ge4$ budgets within
each task, then weights tasks equally. A seed-mean environment--budget cell is
descriptively collapsed below score $20$. Across proximal $K=1,2,3,4$, the mean
rises $25.38\to41.20\to51.92\to63.97$ while collapsed cells fall
$34\to22\to16\to8$ of 63. $K=4$ gains $12.06$ over $K=3$ with task-resampling
interval $[3.84,21.11]$ and seven of nine positive task means, yet its cellwise
median is only $+.34$: eight rescues, no reversals, contribute $54.7\%$ of the
gain. The $T\ge10$ means are $20.27,31.01,42.23,56.79$. PART-Lin $K=3$ raises
the high-budget mean from $38.43$ to $46.44$ and reduces collapsed cells from
25 to 21; its $+1.37$ cellwise median indicates a broader shift. Multi-hop thus
shifts the transition chiefly by rescuing the lower tail, not by uniformly
lifting every cell.

\begin{figure*}[t]
\centering
\includegraphics[width=\textwidth]{figures/stability_envelope.pdf}
\caption{Fixed-budget sweep over nine locomotion tasks. (a) Mean score over
tasks. (b) Collapse rates for $T\ge4$, reported both by environment--budget
cell and by seeded run. All curves and cell means use four seeds. Both realizations broaden the stable
high-budget region with depth.}
\label{fig:stability}
\end{figure*}

\begin{table}[t]
\caption{Aggregate fixed-budget results. Every entry uses four seeds. Best
($T$) is the descriptive maximum over the shared 14-value grid; High denotes
$T\ge4$. Collapse entries are seed-mean cells / raw seeded runs.}
\centering
\small
\setlength{\tabcolsep}{3pt}
\begin{tabular}{lrrrr}
\toprule
Method & $T=1.5$ & Best ($T$) & High & Collapse\\
\midrule
TD3+BC & 74.79 & 74.79 (1.5) & 25.38 & 34 / 142\\
PART-P2 & 78.05 & 80.66 (2.5) & 41.20 & 22 / 111\\
PART-P3 & \textbf{78.95} & 81.41 (2.5) & 51.92 & 16 / 84\\
PART-P4 & 77.70 & \textbf{82.06 (4)} & \textbf{63.97} & \textbf{8 / 54}\\
PART-L2 & 75.13 & 78.43 (2.5) & 38.43 & 25 / 112\\
PART-L3 & 77.74 & 78.20 (2.5) & 46.44 & 21 / 95\\
\bottomrule
\end{tabular}
\label{tab:stability}
\end{table}

The high-budget aggregate is a robustness estimand, not a deployment rule.
At TD3+BC's grid-best shared budget, $T=1.5$, PART-P3 and P4 score $78.95$
and $77.70$ versus $74.79$. Allowing one globally shared grid value per
method, PART-P4 reaches $82.06$ at $T=4$. In a post-hoc tuning-transfer
sensitivity that selects each method's $T$ on two dynamics families and tests
on the third, P4 averages $78.44$ versus $74.79$ for TD3+BC ($+3.65$; 7/9
task differences positive). Family gains vary, so these summaries demonstrate
a practical regime but do not provide an offline selection rule.

\paragraph{Task dependence and replication.}
Responses range from smooth to sharply nonmonotone, so there is no universal
critic-valid radius. Nominal $T$ is the strongest absolute within-environment
rank predictor of return in eight of nine environments; realized movement is
a mechanism readout, not a substitute for the intervention. P4--P3 is positive
within all three dynamics families ($+.79,+13.04,+22.33$), and every
leave-one-family-out mean remains positive ($+6.92$ to $+17.69$). Both
host-separated seed pairs reproduce the high-budget ordering:
$25.97,42.02,52.78,63.61$ and $24.79,40.38,51.05,64.33$. These are disjoint-seed
replications of one implementation.

The linearized sweep shows that the gain is not unique to the proximal loss.
At high budgets, PART-Lin with $K=3$ gains $8.01$ points over $K=2$
(95\% task-resampling interval $[3.56,13.53]$; eight of nine positive task means)
and reduces collapse cells from $25$ to $21$. Both disjoint seed pairs reproduce
the direction, with gains of $7.23$ and $8.80$. PART-Prox remains stronger overall, so these
are comparisons of complete training procedures rather than isolated
integrator effects.

\subsection{Transition and large-budget failure are nonlocal}


At $T=20$, the projected PART-Lin $K=2$ target is clipped at an action bound in
$94$--$95\%$ of states, and the effective action-step coefficient
$h_{\rm eff}$ is only $14\%$ and $9\%$ of the nominal $dT/K$ coefficient in
the two hops. These are medians over 18 environment--seed runs per hop. The
practical operation is projected target generation followed by amortized neural
regression under a moving critic. The sweep measures the resulting end-to-end
stability envelope; local truncation error lies outside this estimand.

Figure~\ref{fig:failure} shows the common failure signature. Environment
return supplies the external label: collapsed runs score below $20$. These
runs have orders-of-magnitude larger TD-error tails, larger critic magnitudes,
and larger realized movements. The controlled common critic is the frozen
final TD3+BC $T=1$ critic matched by environment and seed. A run-total gain
sums the mean $Q$ change across actor hops on the shared state batch. For
collapsed PART-Lin runs, its median under the jointly trained critic is
$3.33\times10^{11}$; under the common critic it is $-1.42$, compared with
$+2.11$ for stable runs. For PART-Prox, the common-critic gain changes from
$+2.51$ in stable runs to $-1.06$ in collapsed runs. This sign reversal
diagnoses cross-critic disagreement consistent with
Eq.~\eqref{eq:critic_error_margin}. A separate four-environment seed-2/3
sensitivity audit,
using a fixed TD3+BC $T=.05$ reference critic, preserves the split: all 45
stable gains are positive and all 35 collapsed gains are negative (Supplement).

A simulator-reference audit provides targeted evidence of critic miscalibration in the screened checkpoints: the
median signed error $\widehat Q_1(s,a_K)-Q_{\rm MC}^{\rho_1}(s,a_K)$ is
$+1.70\times10^{12}$ in five collapsed checkpoints and $-41.4$ in 55 stable
ones, using the median of checkpoint-level state medians. Episode-cluster
intervals exclude zero for critic error at 58 of 60 checkpoints, and restoring
the sampled episode clock preserves the error sign at all 60. Its action-ranking
subset is sparse, so we use it only as targeted checkpoint-level evidence.

\begin{figure*}[t]
\centering
\includegraphics[width=\textwidth]{figures/failure_signature.pdf}
\caption{Median diagnostics over the matched
$T\in\{4,7,10,14,20\}$ audit. Proximal and linearized realizations share a
diagnostic pattern consistent with critic miscalibration: critic errors and
realized movements increase sharply, and run-total gains under a common critic become
negative. Environment return labels failure; the common critic diagnoses
cross-critic disagreement.}
\label{fig:failure}
\end{figure*}

\subsection{How multi-hop helps: exposure versus reach}

The phase diagram points to an exposure--reach mechanism rather than a
small-step accuracy effect. For fixed $T$, only PART's first $T/K$ branch enters
the Bellman target, while the later branches can extend the deployed policy.

On identical next-state batches, P3/P4 first-actor displacement is below
TD3+BC in 88/90 and 89/90 cells (median squared-displacement ratios
$.650,.565$), while final-actor displacement is below it in only 37/90 and
44/90 ($1.060,1.025$). In every P3 and P4 cell the final actor lies farther
from the paired next dataset action than the first Polyak actor (median RMS
ratios $1.30,1.36$). Under one matched
TD3+BC target critic and identical smoothing,
the operational TD-target perturbation is also lower in 86/90 and 88/90 cells
($.717,.617$), with all nine task medians below one. These learned-critic
readouts establish asymmetric movement and reduced critic-space perturbation,
not critic accuracy or return causality. Seed-2/3 and lower-exposure return
counterexamples are reported in the Supplement.

Bootstrap exposure denotes this target-action-displacement proxy, not distance
to dataset support, critic error, or Bellman-target error.

PART-Prox also realizes a smaller first-hop proxy than PART-Lin once $T\ge7$.
Exactly one realization collapses in 11/90 matched cells; in every case the
stable run has smaller first and final movement. This associates controlled
movement with stability without isolating a causal route.

\subsection{Separation explains part of the multi-hop gain}

We test sequential re-centering with a two-seed control whose directly
data-anchored target/deployed actors use $T/3$ and $T$, compared only with a
contemporaneous PART-P3 rerun.
\begin{table}[t]
\caption{P3 policy-separation control over 90 matched task--budget--seed pairs.
Collapse is the descriptive score-20 count before / after seed averaging.}
\centering
\small
\setlength{\tabcolsep}{2.5pt}
\begin{tabular}{lrrr}
\toprule
Condition & Mean & Raw $<20$ & Cell $<20$\\
\midrule
PART-P3 rerun &57.67&23/90&9/45\\
Two-actor control &59.22&21/90&9/45\\
\bottomrule
\end{tabular}
\label{tab:policy_separation}
\end{table}
Control minus PART is $+1.56$ with task-resampling interval $[-1.52,5.44]$,
while its paired median is $-.42$ and PART wins 52/90 pairs; only 2/9 task
means favor the control. Thus policy separation attains the observed P3
aggregate regime without sequential re-centering. This is a bundled,
non-compute-matched comparison that also changes anchoring and critic
trajectory, not an equivalence or re-centering-only effect. The completed P4
control uses directly data-anchored actors with coefficients $T/4$ and $T$.
Across 90 matched pairs, PART-P4 minus control has task-equal mean $-2.31$,
task-resampling interval $[-7.25,2.41]$, paired median $+.68$, and
wins/ties/losses $55/1/34$. The predeclared outcome is \emph{unresolved},
not equivalence or superiority. Every pair is within one host stack; the two
45-pair host blocks use different resolved stacks, which we disclose without
discarding the author-verified finite scores (Supplement).

Eight-run controls give $71.4,70.0,59.4$ for compute-matched sequential,
fixed-reference, and direct-final refinement; slower target averaging does not
rescue $K=1$. Route shadows miss their common-continuation rule (Supplement).

\subsection{Regime synthesis and limitations}

The results support one regime-based reading. At small horizons, the
Wasserstein/JKO view explains the common local direction, and ReLU residence
makes the exact affine identity empirically relevant. It does not predict a
return ordering. As $T$ leaves the small-budget plateau, TD3+BC couples a large
policy move directly into bootstrapping; PART reduces that first-branch exposure
while retaining final reach, and the stability advantage opens. In the extreme
tail, all procedures become nonlocal and degrade, but added depth moves the
collapse boundary and rescues its lower tail. Policy separation is sufficient
for the P3 aggregate, while the P4 chain increment remains unresolved. The live
$K$ intervention still bundles step size, re-centering, compute, normalization,
routing, and optimization history.

Both headline score grids use four seeds, and the proximal grid forms two
host-separated pairs. The full P3 movement and exposure audit uses seeds 0--1;
a restricted four-environment seed-2/3 exposure audit supplies a sensitivity
check, while the P4 target dump uses four seeds. For proximal seeds 2--3,
provenance consists of the archived launch configuration and score matrices;
exact per-run manifests were not retained, and post-hoc P4 manifests do not
reconstruct them. The protocol-v2 full P4 final-displacement audit is limited
to seeds 0--1. The tasks span three dynamics families, collapse is descriptive, and controls
are post hoc. The P3 and P4 two-actor controls use seeds 0--1 and are not compute matched.
The P3 control includes ten Walker2d-expert cells completed through a
documented CPU recovery chain (eight resumed and two CPU-only); the P4
comparison forms two 45-pair host blocks with different disclosed runtime
stacks. We provide no
offline rule for choosing $T$ or $K$. PART makes $K$ actor optimizer calls per
delayed update and stores $K$ persistent actors, so actor work and persistent
actor state scale as $O(K)$. On one H200/JAX 0.10.2 stack, compiled actor-phase
time rises from $1.435$ to $1.887$ ms and scoped backend high-water memory from
9.68 to 18.05 MB between $K=1$ and $K=4$ (1.316$\times$, 1.863$\times$); these
one-stack measurements are not end-to-end scaling. The ideal fixed-critic
analysis motivates but does not certify the live chain.

\section{CONCLUSION}
\label{sec:conclusion}

PART turns the JKO motivation into a finite-step stability intervention. The
small-$T$ limit explains why ideal refinements share a local flow, but the
empirical advantage appears when TD3+BC leaves that regime. Across four seeds,
the $T\ge4$ mean rises from $25.38$ at $K=1$ to $63.97$ at $K=4$ as low-return
cells fall from $34/63$ to $8/63$; both disjoint seed pairs reproduce the
ordering. Matched audits reveal the corresponding signature: added depth
contracts the target-coupled first branch and its critic-space perturbation
without typically contracting the deployed policy, while most of the P4 gain
comes from lower-tail rescue. At extreme budgets the local theory is
inapplicable and every method degrades, so depth shifts rather than removes the
failure boundary. The P3 control identifies policy separation as a simpler
sufficient mechanism hypothesis, while the completed P4 comparison is
unresolved. The data therefore establish PART's procedure-level robustness in
the mid/large-$T$ stress region, but not an incremental return benefit from
sequential re-centering itself.

\bibliography{references}
\section*{Checklist}

\begin{enumerate}

  \item For all models and algorithms presented, check if you include:
  \begin{enumerate}
    \item A clear description of the mathematical setting, assumptions, algorithm, and/or model. [Yes/No/Not Applicable]\\
    \textbf{Yes.} Sections~\ref{sec:background}--\ref{sec:method} define the setting, actor chain, losses, and assumptions.
    \item An analysis of the properties and complexity (time, space, sample size) of any algorithm. [Yes/No/Not Applicable]\\
    \textbf{Yes.} Section~\ref{sec:method} states that BAR uses $K$ persistent actors and $K$ actor optimizer calls per delayed update; Eq.~\eqref{eq:margin}, Proposition~\ref{prop:exposure}, and the local-idealization subsection state the theoretical properties and scope.
    \item (Optional) Anonymized source code, with specification of all dependencies, including external libraries. [Yes/No/Not Applicable]\\
    \textbf{Partial.} The anonymized release contains the BAR launch, evaluation, diagnostic, and verification code plus the compact policy-separation audit. The control-run configs identify local training-source files by SHA-256, but did not record a clean Git revision.
  \end{enumerate}

  \item For any theoretical claim, check if you include:
  \begin{enumerate}
    \item Statements of the full set of assumptions of all theoretical results. [Yes/No/Not Applicable]\\
    \textbf{Yes.} Each proposition states its assumptions, and the local-flow assumptions are separated from the coupled large-budget regime.
    \item Complete proofs of all theoretical results. [Yes/No/Not Applicable]\\
    \textbf{Yes.} Complete proofs are provided in the supplementary material.
    \item Clear explanations of any assumptions. [Yes/No/Not Applicable]\\
    \textbf{Yes.} Sections~\ref{sec:method} and~\ref{sec:experiments} explain the learned-critic, fixed-$K$, small-step, and implementation restrictions.
  \end{enumerate}

  \item For all figures and tables that present empirical results, check if you include:
  \begin{enumerate}
    \item The code, data, and instructions needed to reproduce the main experimental results (either in the supplemental material or as a URL). [Yes/No/Not Applicable]\\
    \textbf{Yes.} Released CSVs, figure-generation code, a numerical claim verifier, and build instructions reproduce every main figure and aggregate table.
    \item All the training details (e.g., data splits, hyperparameters, how they were chosen). [Yes/No/Not Applicable]\\
    \textbf{No.} All recoverable settings are listed, but exact per-run BAR-P4 configurations and manifests for proximal seeds 2--3 are unavailable; the policy-separation configs also omit a source revision and environment lock.
    \item A clear definition of the specific measure or statistics and error bars (e.g., with respect to the random seed after running experiments multiple times). [Yes/No/Not Applicable]\\
    \textbf{Yes.} Captions and text define the equally task-and-budget-weighted high-budget estimand, seed-mean cells, raw seeded runs, collapse thresholds, paired ratios, and task-resampling intervals. Both headline grids use four seeds; proximal seeds form two host-separated pairs. Mechanism audits use the explicitly reported two- or four-seed subsets.
    \item A description of the computing infrastructure used. (e.g., type of GPUs, internal cluster, or cloud provider). [Yes/No/Not Applicable]\\
    \textbf{No.} The released artifact records software and execution protocols, including the policy-separation CPU recovery chain, but does not contain a complete, independently verified inventory of all external-machine hardware.
  \end{enumerate}

  \item If you are using existing assets (e.g., code, data, models) or curating/releasing new assets, check if you include:
  \begin{enumerate}
    \item Citations of the creator If your work uses existing assets. [Yes/No/Not Applicable]\\
    \textbf{Yes.} D4RL, TD3+BC, and the CORL implementation basis are cited.
    \item The license information of the assets, if applicable. [Yes/No/Not Applicable]\\
    \textbf{Yes.} The code release identifies its license; upstream dependencies retain their own licenses.
    \item New assets either in the supplemental material or as a URL, if applicable. [Yes/No/Not Applicable]\\
    \textbf{Yes.} An anonymized code and compact-results artifact accompanies the submission.
    \item Information about consent from data providers/curators. [Yes/No/Not Applicable]\\
    \textbf{Not Applicable.} We use public simulated-control benchmark datasets and collect no human data.
    \item Discussion of sensible content if applicable, e.g., personally identifiable information or offensive content. [Yes/No/Not Applicable]\\
    \textbf{Not Applicable.} The assets contain simulated locomotion trajectories and no personal or sensitive content.
  \end{enumerate}

  \item If you used crowdsourcing or conducted research with human subjects, check if you include:
  \begin{enumerate}
    \item The full text of instructions given to participants and screenshots. [Yes/No/Not Applicable]\\
    \textbf{Not Applicable.}
    \item Descriptions of potential participant risks, with links to Institutional Review Board (IRB) approvals if applicable. [Yes/No/Not Applicable]\\
    \textbf{Not Applicable.}
    \item The estimated hourly wage paid to participants and the total amount spent on participant compensation. [Yes/No/Not Applicable]\\
    \textbf{Not Applicable.}
  \end{enumerate}

\end{enumerate}

\end{document}
